\documentclass[twocolumn]{aastex631}
\usepackage{amsmath}
\usepackage{booktabs}
\shorttitle{SDSS mass transition in quenching and optical AGN incidence}
\shortauthors{Feedback transition mass}
\begin{document}

\title{SDSS mass transition in quenching and optical AGN incidence: selection-aware SDSS optical proxy integration}
\author{NebulaMind Research Autopilot}
\affiliation{Public SDSS DR17 data only}

\begin{abstract}
We use the cached SDSS DR17 emission-line subset to identify the stellar-mass regime where quenched fraction and optical AGN incidence rise together. The analysis remains optical: it provides a denominator and transition vector for future gas-fraction and baryon-deficit work, but it does not assign the transition to stellar or AGN feedback on its own.
\end{abstract}

\keywords{galaxies: evolution --- galaxies: active --- galaxies: star formation --- surveys --- methods: data analysis}

\section{Introduction}\label{sec:purpose}
This analysis preserves the active proposal title, 'Locating the transition from stellar-feedback to AGN-feedback regulation', but narrows the manuscript to the directly measured SDSS optical quantities reported below. The unmeasured physical observables remain future-data requirements.

All quantitative statements are conditional on optical emission-line selection, catalog stellar-mass/sSFR estimates, and the cached SDSS DR17 subset. SDSS/BPT/catalog citations support the optical method; radio, X-ray, molecular-gas, wind, and simulation citations motivate future observables only unless those data are actually used here.


\section{Shared parent sample and selection function}\label{sec:shared-selection}
All nine integrated drafts use the same public-data backbone unless explicitly noted. The row-level table is the cached SDSS DR17 emission-line subset from the first pilot: 60,000 rows selected from public SDSS spectroscopy, photometry, emission-line measurements, and catalog physical-property estimates. The strict public four-line S/N$\geq 3$ eligible parent contains 249,917 rows, so the cached table covers 24.0\% of that strict parent. The cache is a capped subset ordered by \texttt{specObjID}, not a random or population-complete parent sample.

\begin{deluxetable*}{lrrr}
\tabletypesize{\scriptsize}
\tablecaption{Shared SDSS DR17 selection cascade used before paper-specific quantities.\label{tab:selection-cascade}}
\tablehead{\colhead{Selection stage} & \colhead{Public DR17 rows} & \colhead{Cached rows} & \colhead{Retention vs. spectro-z parent (fraction)}}
\startdata
SpecObj GALAXY, 0.02<z<0.12 & 501,060 & -- & 1.000 \\
plus galSpecInfo/PhotoObj/galSpecExtra and mass/sSFR bounds & 416,554 & -- & 0.831 \\
plus galSpecLine join & 416,554 & -- & 0.831 \\
four BPT lines positive with positive errors & 373,445 & 60,000 & 0.745 \\
four BPT lines S/N$\geq 3$ & 249,917 & 60,000 & 0.499 \\
four BPT lines S/N$\geq 5$ & 176,523 & 42,446 & 0.352 \\
four BPT lines S/N$\geq 10$ & 91,768 & 22,311 & 0.183 \\
\enddata
\tablecomments{Counts are read-only public SDSS DR17 count queries plus the cached local CSV. Cached rows are shown only where the cache applies.}
\end{deluxetable*}

The four-line requirement is strongly selection dependent. In the public counts, S/N$\geq 3$ keeps 33.6\% of the $-12<\log {\rm sSFR}<-11$ parent bin but 94.9\% of the $-10<\log {\rm sSFR}<-9.5$ bin. Therefore every incidence, quenching, density, gas-denominator, or target-vector statement in this integration is conditional on the four-line emission-line selection.

Cached-versus-public marginal checks found no redshift, stellar-mass, or sSFR bin with a cached-minus-public fraction difference above 5 percentage points. The largest absolute differences were 2.03 percentage points in redshift, -1.63 percentage points in stellar mass, and -0.58 percentage points in sSFR. This is a representativeness diagnostic only; it does not make the capped cache random or complete.


\section{Measurements}\label{sec:measurements}
The row-level measurements include redshift, stellar mass, catalog specific star-formation rate, model colors, and four optical line fluxes/errors. BPT labels and all derived denominators are recomputed locally from the cached table. Catalog SFR/sSFR values are treated as low-redshift SDSS physical-property estimates rather than direct resolved gas or feedback measurements \citep{sdssdr17,brinchmann2004,york2000}.


\section{Optical denominator for feedback-transition mass}\label{sec:topic-result}
The consolidated proposal question is: At what stellar-mass scale do quenched fraction and optical AGN incidence rise in the same SDSS denominator? The integrated manuscript deliberately demotes the claim to the directly measured SDSS quantity. It is not a full physical-feedback test.

\begin{itemize}
\item The first stellar-mass bin with quenched fraction above 0.5 is 11.0-12.5.
\item The optical AGN fraction peaks in the 11.0-12.5 bin at 0.520.
\item The result is an optical transition diagnostic; gas fractions and baryon deficits are needed before assigning the transition to stellar or AGN feedback.
\end{itemize}


\begin{figure}
\centering
\includegraphics[width=\columnwidth]{../figures/fig-topic.pdf}
\caption{SDSS DR17 optical denominator/proxy diagnostic for the feedback-transition mass vector. The figure summarizes the cached optical result used for target definition.}
\label{fig:topic}
\end{figure}

\section{Interpretation and missing observables}\label{sec:missing}
SDSS-only pilot; full proposal requires the additional survey data named in the research-topic page. The full proposal requires: gas fractions, baryon deficits, halo masses, stellar-feedback observables, and high-redshift extensions.

Mass, color bimodality, halo shock, central/satellite, and black-hole-mass studies define variables that must be added before attributing a mass vector to a physical feedback transition \citep{kauffmann2003mass,baldry2004,peng2010,peng2012,dekel2006,bluck2023,piotrowska2022}.


\section{Reproducibility and safety}\label{sec:repro}
This manuscript uses the cached SDSS DR17 subset and the accompanying local manifest of inputs and outputs. The output is a local draft PDF and manifest entry only. No public-linked PDF was replaced.

\section{Conclusion}\label{sec:conclusion}
The first stellar-mass bin with quenched fraction above 0.5 is 11.0-12.5, and the optical AGN fraction peaks at 0.520 in that same bin. These values define an optical transition vector, but gas fractions, baryon deficits, and halo-scale measurements are still needed before a causal feedback interpretation.


\begin{thebibliography}{}
\bibitem[Abdurro'uf et al.(2022)]{sdssdr17} Abdurro'uf, Accetta, K., Aerts, C., et al. 2022, ApJS, 259, 35
\bibitem[Brinchmann et al.(2004)]{brinchmann2004} Brinchmann, J., Charlot, S., White, S.~D.~M., et al. 2004, MNRAS, 351, 1151
\bibitem[Kauffmann et al.(2003b)]{kauffmann2003mass} Kauffmann, G., Heckman, T.~M., White, S.~D.~M., et al. 2003b, MNRAS, 341, 33
\bibitem[York et al.(2000)]{york2000} York, D.~G., Adelman, J., Anderson, J.~E., Jr., et al. 2000, AJ, 120, 1579

\bibitem[Baldry et al.(2004)]{baldry2004} Baldry, I.~K., Glazebrook, K., Brinkmann, J., et al. 2004, ApJ, 600, 681
\bibitem[Bluck et al.(2023)]{bluck2023} Bluck, A.~F.~L., Piotrowska, J.~M., \& Maiolino, R. 2023, ApJ, 944, 108
\bibitem[Dekel \& Birnboim(2006)]{dekel2006} Dekel, A., \& Birnboim, Y. 2006, MNRAS, 368, 2
\bibitem[Peng et al.(2010)]{peng2010} Peng, Y.-j., Lilly, S.~J., Kovac, K., et al. 2010, ApJ, 721, 193
\bibitem[Peng et al.(2012)]{peng2012} Peng, Y.-j., Lilly, S.~J., Renzini, A., \& Carollo, M. 2012, ApJ, 757, 4
\bibitem[Piotrowska et al.(2022)]{piotrowska2022} Piotrowska, J.~M., Bluck, A.~F.~L., Maiolino, R., \& Peng, Y.-j. 2022, MNRAS, 512, 1052
\end{thebibliography}

\end{document}
